% % ---
% % Capitulo de Resultados
% % ---

\chapter{Resultados}\label{cap:resultados}

A proposta deste software é auxiliar os terapeutas na implementação do método ABA, oferecendo uma gama de recursos para coleta de dados e anotações facilitadas, além de facilitar o compartilhamento com outros terapeutas da criança e a família, aumentando a eficiência no acesso aos dados, tanto brutos quanto compilados. O software será desenvolvido com foco na usabilidade e acessibilidade para atingir um público mais amplo.

A interface será construída utilizando Flutter Dart, consumindo uma API baseada em Quarkus JAVA. O software terá como uma de suas principais funcionalidades a capacidade de criar e consultar anotações voltadas para a aplicação do método ABA. Além disso, contará com uma seção destinada especificamente à técnica do ABA, treino por tentativas discretas.

\section{Levantamento Bibliográfico}

Com o intuito de explorar as dificuldades enfrentadas pelos terapeutas ao aplicar o método ABA, foi realizada uma pesquisa bibliográfica. O objetivo era compreender conceitos, técnicas e soluções para desafios associados a esse método. Com base nessa pesquisa, foi delineada uma abordagem para a criação de anotações compartilhadas e compiladas, visando facilitar a compreensão do progresso e das dificuldades individuais de cada paciente.

Adicionalmente, observou-se que uma das técnicas do ABA utiliza estímulos específicos, e dependendo da resposta do paciente, são feitas anotações distintas. Com isso em mente, foi concebido um método para automatizar o armazenamento desses dados, permitindo um registro mais eficiente e sistemático. Detalhes da pesquisa bibliográfica podem ser encontrados no capitulo \ref{cap:referencial_teorico}

\section{Levantamento de Funcionalidades}

As funcionalidades propostas se concentram primariamente no levantamento, compartilhamento e análise de dados, oferecendo amplas opções de configuração para tornar o software altamente flexível e adaptável às preferências do usuário. Além disso, incluem a capacidade de desvincular um paciente do terapeuta ou da rede de terapeutas.

O levantamento de dados pode ser realizado de três maneiras distintas: por meio de texto escrito em linguagem natural, com a inserção em tabelas configuráveis com opções predefinidas, ou pela aplicação de uma técnica do ABA, envolvendo interações do paciente com cartões, registrando dados como sucesso, sucesso com dica e erro.

Os dados podem ser compartilhados entre os terapeutas na rede e a família do paciente. A cada anotação realizada pelo terapeuta, é dada a opção de decidir se esses dados serão compartilhados com a família do paciente. Além disso, há a possibilidade de o terapeuta bloquear temporariamente ou permanentemente a capacidade da família de fazer anotações adicionais.

A análise dos dados pode ser realizada de maneira geral, ou de forma detalhada, acessando cada anotação individualmente. Enquanto o terapeuta possui acesso às informações de todos os pacientes, a família só pode visualizar os dados referentes ao paciente vinculado.

A desvinculação do paciente do terapeuta pode ser executada por qualquer usuário após a confirmação por senha, e todos os outros vinculados ao paciente serão informados. Após a desvinculação, o terapeuta perde o acesso aos dados do paciente, enquanto a família mantém acesso a todas as anotações relacionadas ao paciente.

\section{Proposta do layout}

Para o desenvolvimento do layout foi preciso uma análise de publico alvo, 
para um melhor entendimento foi criado duas personas para que o desenvolvimento do layout seja mais adequado com o publico alvo. 
O layout pode ser visto na seção \ref{sec:resultados_esperados}. 

\section{Simulação de Uso com Dados Fictícios}
